\begin{figure}[htbp]
\centering
\begin{tikzpicture}[x=.44cm,y=.44cm,font=\small,
    loop/.style={draw=ymInk,line width=1.05pt,line join=round},
    direction/.style={-{Stealth[length=1.65mm,width=1.35mm]},ymInk}]
  \fill[ymInk!9] (0,0) rectangle (5,5);
  \foreach \x in {0,...,20} \draw[ymGrid,very thin] (\x,0) -- (\x,10);
  \foreach \y in {0,...,10} \draw[ymGrid,very thin] (0,\y) -- (20,\y);
  \foreach \x in {0,5,...,20}
    \draw[black!62,densely dashed] (\x,0) -- (\x,10);
  \foreach \y in {0,5,10}
    \draw[black!62,densely dashed] (0,\y) -- (20,\y);
  % Each adjacent cell is a reflection of the preceding cell.
  \foreach \tx/\sx in {0/1,10/-1,10/1,20/-1}{
    \foreach \ty/\sy in {0/1,10/-1}{
      \begin{scope}[shift={(\tx,\ty)},xscale=\sx,yscale=\sy]
        \draw[loop] (1,1) -- (3,1) -- (3,2) -- (1,2) -- cycle;
        \draw[direction] (1.35,1) -- (2.42,1);
        \fill[ymInk] (1,1) circle[radius=1.6pt];
      \end{scope}
    }
  }
  \draw[black!65,|-|] (0,-.46) -- (5,-.46)
    node[midway,fill=white,inner xsep=3pt] {$\ell$};
\end{tikzpicture}
\caption[Reflected loops in the chessboard estimate]{Part of a coordinate
section of the chessboard tiling. Starting with the based loop in the shaded cell,
reflect its links across each dashed seam. One reflection reverses the
orientation in this plane and two perpendicular reflections restore it.
Every copy lies strictly inside its cell. The product in
\eqref{ch01:chessboard-root} uses the copies in all four directions.}
\label{ch01:fig-chessboard}
\end{figure}
